\section{El cultivo}\label{sec:cultivo}
Las parcelas están en los Valles Altos de la Mesa Central, una de las principales regiones productoras de cebada del
país. Ahí la cebada es \emph{de temporal}, es decir, sin riego, y toda el agua viene de la lluvia del verano. Las
parcelas están entre unos 2{,}200 y 2{,}700 metros de altitud.

\subsection{Calendario}
El ciclo dura de tres a cuatro meses (Tabla~\ref{tab:calendario}). La siembra depende de cuándo llegan las lluvias,
así que dos parcelas pueden estar en etapas distintas el mismo día. Por eso conviene comparar las parcelas según su
etapa y no según la fecha del calendario.

\begin{table}[!t]
\caption{Calendario aproximado del cultivo}\label{tab:calendario}
\centering\footnotesize
\begin{tabularx}{\columnwidth}{@{}>{\raggedright\arraybackslash}p{2.3cm}>{\raggedright\arraybackslash}p{1.7cm}X@{}}
\toprule
Etapa & Fecha aproximada & Relevancia \\
\midrule
Preparación del terreno & Abril y mayo & Visible en los índices de labranza y residuos \\
Siembra & Mayo a mediados de junio & Depende de la llegada de las lluvias \\
Amacollamiento & 10 a 35 días después & Define cuántos tallos tendrá la planta \\
Encañe y embuche & Junio y julio & Inicio del periodo crítico \\
Espigamiento y floración & Julio y agosto & Máximo verdor de la parcela \\
Llenado de grano & Agosto y septiembre & Define el peso del grano \\
Madurez y cosecha & Septiembre y octubre & El verdor cae rápido \\
\bottomrule
\end{tabularx}
\end{table}

\subsection{Periodo crítico}
El rendimiento de la cebada depende sobre todo del número de granos por metro cuadrado. Ese número se define en una
ventana corta, el \emph{periodo crítico}, que va de unas tres o cuatro semanas antes de la floración hasta unos días
después. En cebada esa ventana ocurre sobre todo antes del espigamiento~\cite{arisnabarreta2008}. El exceso de agua
también tiene su periodo crítico, en las semanas previas a la floración~\cite{desanceledonio2014}.

\subsection{Información que no tenemos}
El conjunto entregado no incluye la variedad, la fecha de siembra, la fertilización ni el manejo de cada productor.
La fecha de siembra se puede estimar desde el satélite, a partir del momento en que empieza a crecer el verdor. El
resto queda como error que ningún modelo podrá explicar.
